% APA 7th edition manuscript template (LaTeX)
% Requires: apa7 + biblatex + biblatex-apa + biber (Overleaf supports this).
%
% Compile (local): latexmk -pdf -interaction=nonstopmode -bibtex -usebiber apa7-manuscript.tex
% Compile (Overleaf): set "Biber" as the bibliography tool.

\documentclass[man]{apa7}

\usepackage[american]{babel}
\usepackage{csquotes}
\usepackage{amsmath,amssymb}
\usepackage{graphicx}
\usepackage{booktabs}

\usepackage[style=apa,backend=biber,natbib=true,sortcites=true]{biblatex}
\addbibresource{references.bib}

\title{Your Paper Title: A Subtitle if Needed}
\shorttitle{SHORT TITLE}

% --- Authors and Affiliations ---
\authorsnames[1,2,1]{First A. Author, Second B. Author, Third C. Author}
\authorsaffiliations{
  {Department of Psychology, University of Example},
  {School of Education, Another University}
}

% --- Author Note ---
\authornote{
  \addORCIDlink{First A. Author}{0000-0000-0000-0000}

  Correspondence concerning this article should be addressed to
  First A. Author, Department of Psychology, University of Example,
  123 Academic Way, City, State 12345.
  E-mail: first.author@example.edu
}

\abstract{Write a 150--250 word abstract that summarizes the problem,
method, main results, and implications. The abstract should be a single
paragraph without citations. It should state the purpose, describe the
participants and methods, report the main findings with effect sizes,
and note the conclusions and implications.}

\keywords{keyword1, keyword2, keyword3, keyword4, keyword5}

\begin{document}
\maketitle

% ============================================================
\section{Introduction}
% ============================================================

Introduce the research problem and its significance. Provide
background that motivates the study \parencite{smith2023}. State
the research question or hypothesis clearly.

\subsection{Theoretical Framework}

Describe the theoretical lens through which this study is framed.
\textcite{jones2022} proposed a framework that is relevant to the
current investigation.

\subsection{Present Study}

State the purpose of the current study and the specific hypotheses.
Based on prior work \parencite{smith2023,jones2022}, we hypothesize
that\ldots

% ============================================================
\section{Literature Review}
% ============================================================

\subsection{Theme 1}

Summarize the relevant literature organized by theme rather than
chronologically. \textcite{author2020} established foundational
findings in this area.

\subsection{Theme 2}

Discuss converging and diverging evidence. Several studies have
examined this phenomenon \parencite{other2021,diff2022}, with
mixed results.

% ============================================================
\section{Method}
% ============================================================

\subsection{Participants}

Describe the sample. Participants were $N = 120$ undergraduate
students ($M_{\text{age}} = 19.8$, $SD = 1.4$; 65\% female).

\subsection{Materials}

Describe instruments, measures, and their psychometric properties.

\subsection{Procedure}

Describe the study procedure in sufficient detail for replication.

\subsection{Data Analysis}

Describe the analytic strategy, including any software used.

% ============================================================
\section{Results}
% ============================================================

Report results in order of the hypotheses. Use tables and figures
to present key findings.

% --- Example Table ---
\begin{table}[ht]
  \caption{Descriptive Statistics by Condition}
  \label{tab:descriptives}
  \begin{tabular}{lrrr}
    \toprule
    Condition & $n$ & $M$ & $SD$ \\
    \midrule
    Control       & 60 & 4.21 & 1.03 \\
    Experimental  & 60 & 5.47 & 0.98 \\
    \bottomrule
  \end{tabular}
\end{table}

Results indicated a significant effect, $t(118) = 3.42$,
$p < .001$, $d = 0.63$, 95\% CI $[0.27, 0.99]$.

% --- Example Figure ---
\begin{figure}[ht]
  \centering
  % \includegraphics[width=0.8\textwidth]{figures/results-plot.pdf}
  \caption{Mean scores by condition. Error bars represent 95\%
  confidence intervals.}
  \label{fig:results}
\end{figure}

% --- Example Equation ---
The regression model is specified as:
\begin{equation}
  Y_i = \beta_0 + \beta_1 X_{1i} + \beta_2 X_{2i} + \varepsilon_i
  \label{eq:regression}
\end{equation}

where $\varepsilon_i \sim N(0, \sigma^2)$.

% ============================================================
\section{Discussion}
% ============================================================

Interpret the findings in relation to the hypotheses and prior
literature. The results are consistent with \textcite{smith2023},
who found similar effects.

\subsection{Limitations}

Acknowledge methodological limitations and their implications
for interpreting the results.

\subsection{Implications and Future Directions}

Discuss theoretical and practical implications. Suggest directions
for future research that address the identified limitations.

% ============================================================
\section{Conclusion}
% ============================================================

Provide a concise summary of the main findings and their
significance.

% ============================================================
\printbibliography

% ============================================================
% Appendices (if needed)
% ============================================================
\appendix

\section{Supplementary Materials}

Include supplementary tables, figures, or instrument descriptions
that support but are not essential to the main text.

\end{document}
